\documentclass[12pt,a4paper]{article}

% ============================================================
% PACKAGES
% ============================================================
\usepackage[utf8]{inputenc}
\usepackage[T5]{fontenc}
\usepackage[vietnamese]{babel}
\usepackage[margin=2.5cm]{geometry}
\usepackage{graphicx}
\usepackage{booktabs}
\usepackage{tabularx}
\usepackage{longtable}
\usepackage{hyperref}
\usepackage{xcolor}
\usepackage{listings}
\usepackage{amsmath}
\usepackage{float}
\usepackage{caption}
\usepackage{subcaption}
\usepackage{enumitem}
\usepackage{titlesec}
\usepackage{fancyhdr}
\usepackage{tikz}
\usetikzlibrary{shapes.geometric, arrows, positioning}

% ============================================================
% THIẾT LẬP TIKZ STYLES
% ============================================================
\tikzstyle{layer} = [rectangle, rounded corners, minimum width=3cm, minimum height=1.2cm, text centered, draw=blue!80, fill=blue!5, text width=2.8cm, font=\small]
\tikzstyle{dropout} = [rectangle, rounded corners, minimum width=3cm, minimum height=1.2cm, text centered, draw=orange!80, fill=orange!5, text width=2.8cm, font=\small]
\tikzstyle{dense} = [rectangle, rounded corners, minimum width=3cm, minimum height=1.2cm, text centered, draw=green!70!black, fill=green!5, text width=2.8cm, font=\small]
\tikzstyle{arrow} = [thick,->,>=stealth]

% ============================================================
% THIẾT LẬP HYPERREF & LISTINGS
% ============================================================
\hypersetup{
    colorlinks=true,
    linkcolor=blue,
    urlcolor=blue,
}

\lstset{
    basicstyle=\ttfamily\small,
    breaklines=true,
    frame=single,
    backgroundcolor=\color{gray!10},
}

\pagestyle{fancy}
\fancyhf{}
\rhead{Nhóm 11 --- Báo cáo Mô hình Học máy}
\lhead{Môn Nhập môn Học máy}
\rfoot{\thepage}
\setlength{\parindent}{0pt}
\setlength{\parskip}{0.6em}

% ============================================================
% BẮT ĐẦU TÀI LIỆU
% ============================================================
\begin{document}

% ── TRANG BÌA ──────────────────────────────────────────────
\begin{titlepage}
    \centering
    \vspace*{1cm}
    {\Large \textbf{ĐẠI HỌC KHOA HỌC TỰ NHIÊN TP.HCM}}\\[0.3cm]
    {\large Khoa Công nghệ Thông tin}\\[2cm]

    {\Huge \textbf{BÁO CÁO MÔ HÌNH HỌC MÁY}}\\[0.5cm]
    {\large Xây dựng và Triển khai Hệ thống Học máy Ứng dụng}\\[2cm]

    {\LARGE \textbf{Dự đoán Chất lượng Không khí Việt Nam}}\\[0.3cm]
    {\large và Cảnh báo Nguy cơ Mắc Bệnh Hô hấp}\\[2cm]

    \begin{tabular}{cll}
        \toprule
        \textbf{STT} & \textbf{MSSV} & \textbf{Họ và tên} \\
        \midrule
        1 & 23120062 & Trần Kim Ngọc \\
        2 & 23120063 & Nguyễn Thành Nguyên \\
        3 & 23120084 & Nguyễn Mạnh Thắng \\
        4 & 23120059 & Trần Đình Luân \\
        5 & 23120064 & Nguyễn Thiện Nhân \\
        \bottomrule
    \end{tabular}\\[2cm]

    {\large Nhóm 11}
\end{titlepage}

\tableofcontents
\vspace{0.5cm}
\noindent\textbf{Repository GitHub:} \href{https://github.com/kimngoc280105/aqi-vietnam}{https://github.com/kimngoc280105/aqi-vietnam}
\newpage

% ============================================================
% 1. GIỚI THIỆU BÀI TOÁN
% ============================================================
\section{Giới thiệu bài toán}

Ô nhiễm môi trường không khí tại Việt Nam, đặc biệt là nồng độ bụi mịn $PM_{2.5}$ tại các đô thị lớn như Hà Nội, TP.HCM và Đà Nẵng đang là vấn đề lớn đe dọa đến sức khỏe cộng đồng. Nhóm hướng đến giải quyết bài toán có ý nghĩa thực tiễn cao: \textbf{Dự báo chuỗi nồng độ bụi mịn $PM_{2.5}$ trong không khí cho 24 giờ tới ($t+1$ đến $t+24$)} và \textbf{Phân loại mức độ nguy cơ sức khỏe cảnh báo sớm}. Việc dự báo trước 24 giờ là cực kỳ quan trọng đối với các nhóm đối tượng nhạy cảm như người mắc bệnh hô hấp mãn tính (hen suyễn, viêm phế quản, COPD) để họ có đủ thời gian chuẩn bị các biện pháp y tế và lập kế hoạch hoạt động an toàn trước một ngày.

Trong các thiết kế ban đầu của hệ thống, chỉ số chất lượng không khí tổng hợp (AQI) được cân nhắc làm biến mục tiêu dự đoán. Tuy nhiên, qua phân tích khám phá dữ liệu (EDA), nhóm đã thực hiện bước chuyển đổi chiến lược (\textit{Target Pivot}) sang dự đoán trực tiếp nồng độ vật chất tuyệt đối $PM_{2.5}$ vì các lý do sau:
\begin{enumerate}
    \item \textbf{Hiện tượng che khuất (Masking effect):} Chỉ số AQI tổng hợp được tính bằng hàm $Max$ nồng độ của nhiều loại khí ($PM_{2.5}, O_3, NO_2, SO_2, CO$). Khi một loại khí giảm đột biến (ví dụ $O_3$ vào ban đêm) và bụi mịn $PM_{2.5}$ tăng cao, AQI tổng thể có thể không phản ánh được mức tăng cục bộ nguy hiểm của bụi mịn.
    \item \textbf{Độ mượt toán học:} $PM_{2.5}$ tuân theo tự tương quan chuỗi thời gian ($Autocorrelation$) liên tục và bám sát các quy luật vật lý (khuếch tán, phát tán) tốt hơn chuỗi AQI vốn mang tính giật cục do các phép làm mượt 24 giờ.
    \item \textbf{Tính thực tế y khoa:} $PM_{2.5}$ là tác nhân nguy cơ chính xâm nhập sâu vào phế nang phổi và ngấm vào máu. Dự đoán trực tiếp $PM_{2.5}$ giúp nâng cao tính tường minh (Explainable AI) và độ tin cậy trong các cảnh báo sức khỏe y tế.
\end{enumerate}

% ============================================================
% 2. TỔNG QUAN DỮ LIỆU ĐẦU VÀO
% ============================================================
\section{Tổng quan dữ liệu đầu vào}

\subsection{Dữ liệu huấn luyện và Tỷ lệ chia dữ liệu}

Dữ liệu thực nghiệm bao gồm 98,187 mẫu thu thập theo giờ từ ngày \textbf{05/08/2022} đến \textbf{29/04/2026} cho ba thành phố Hà Nội, TP.HCM và Đà Nẵng. Vì đặc thù của dữ liệu chuỗi thời gian chứa tính tự tương quan cao, việc chia dữ liệu ngẫu nhiên sẽ dẫn đến hiện tượng rò rỉ thông tin tương lai (\textit{Data Leakage} hay \textit{Look-ahead Bias}). 

Do đó, nhóm sử dụng phương pháp chia dữ liệu theo trình tự thời gian (\textbf{Time-series Split}) với tỷ lệ \textbf{70/15/15}:
\begin{itemize}
    \item \textbf{Tập Huấn luyện (Train) - 70\%}: Dùng để tối ưu hóa trọng số của các mô hình học máy.
    \item \textbf{Tập Kiểm định (Validation) - 15\%}: Dùng để điều chỉnh siêu tham số và kích hoạt các hàm dừng sớm (\textit{Early Stopping}) nhằm chống học tủ.
    \item \textbf{Tập Kiểm thử (Test) - 15\%}: Giữ độc lập hoàn toàn để đánh giá hiệu năng cuối cùng của mô hình trên dữ liệu chưa từng thấy.
\end{itemize}

\subsection{Các bước Tiền xử lý Dữ liệu}

Trước khi đưa dữ liệu vào mô hình, các bước chuẩn bị dữ liệu kỹ thuật sau đã được áp dụng:
\begin{enumerate}
    \item \textbf{Nội suy dữ liệu khuyết thiếu (Imputation):} Sử dụng nội suy tuyến tính (\textit{Linear Interpolation}) kết hợp \textit{Forward Fill} và \textit{Backward Fill} theo từng phân nhóm thành phố độc lập. Tỷ lệ khuyết thiếu thực tế ban đầu cực kỳ thấp ($<0.08\%$) ở các biến trễ thời gian và đạt mức $0\%$ sạch hoàn toàn sau xử lý.
    \item \textbf{Xử lý nhiễu và bất thường (Anomalies):} Phát hiện các giá trị âm ở nồng độ Ozone ($O_3$) do lỗi cảm biến và chuyển thành giá trị trống ($NaN$) để tiến hành nội suy.
    \item \textbf{Giải quyết đa cộng tuyến (Multicollinearity):} Từ ma trận tương quan, biến không gian chỉ số nhà máy trong bán kính 10km ($factories_{10km}$) có tương quan hoàn hảo $r=1.00$ với biến $factories_{5km}$, gây hiện tượng đa cộng tuyến nghiêm trọng. Nhóm quyết định loại bỏ cột $factories_{10km}$, chỉ giữ lại $factories_{2km}$ và $factories_{5km}$.
    \item \textbf{Xử lý ngoại lệ (Outliers):} Áp dụng quy tắc Tukey's Fence mở rộng ($Q3 + 3 \times IQR$). Nhóm quyết định \textbf{giữ lại toàn bộ ngoại lệ} vì chúng phản ánh các đợt ô nhiễm cực đoan có thật trong thực tế (đốt rơm rạ, nghịch nhiệt bức xạ) - những thời điểm mà hệ thống cảnh báo cần hoạt động chính xác nhất.
    \item \textbf{Trích xuất đặc trưng (Feature Engineering):}
    \begin{itemize}
        \item \textit{Đặc trưng trễ (Lag features):} Tạo các cột $PM_{2.5}$ trễ 1h, 3h, 6h, 12h, 24h để nắm bắt xu hướng ngắn hạn.
        \item \textit{Đặc trưng cửa sổ trượt (Rolling features):} Tính trung bình trượt nồng độ $PM_{2.5}$ trong 6h, 24h và 72h để làm mịn nhiễu.
        \item \textit{Mã hóa biến phân loại:} Thực hiện One-Hot Encoding cho các biến địa lý (\texttt{city}) và thời tiết mùa (\texttt{season}).
    \end{itemize}
    \item \textbf{Chuẩn hóa đặc trưng (Scaling):} Sử dụng \texttt{StandardScaler} (Z-score normalization) cho các thuật toán Linear Regression/XGBoost, và \texttt{MinMaxScaler} trong khoảng $[0, 1]$ cho mạng LSTM để hỗ trợ sự hội tụ của hàm kích hoạt phi tuyến. Bộ chuẩn hóa chỉ được khớp ($fit$) trên tập Train và biến đổi ($transform$) trên tập Val/Test nhằm tránh rò rỉ dữ liệu.
\end{enumerate}

% ============================================================
% 3. LỰA CHỌN MÔ HÌNH VÀ KIẾN TRÚC
% ============================================================
\section{Lựa chọn Mô hình \& Kiến trúc}

Nhóm nghiên cứu và phát triển 3 lớp mô hình đại diện cho 3 trường phái tiếp cận:

\begin{enumerate}
    \item \textbf{Mô hình Baseline (SARIMAX):} Lớp mô hình thống kê tuyến tính chuỗi thời gian cổ điển có kết hợp ngoại sinh (nhiệt độ, độ ẩm, tốc độ gió, áp suất). Được dùng làm mốc so sánh cơ sở.
    \item \textbf{Mô hình Học máy nâng cao (XGBoost Regressor):} Thuật toán học máy dựa trên cây quyết định nâng cao cực kỳ mạnh mẽ cho dữ liệu dạng bảng. XGBoost có khả năng học các mối quan hệ phi tuyến tính phức tạp và không nhạy cảm với phân phối lệch của bụi mịn.
    \item \textbf{Mô hình Học sâu tuần tự (LSTM):} Mạng nơ-ron hồi quy có cấu trúc cổng nhớ đặc biệt, thiết kế chuyên biệt để ghi nhớ các phụ thuộc dài hạn trong chuỗi thời gian.
\end{enumerate}

\subsection{Kiến trúc cấu trúc của Mô hình Thống kê SARIMAX}

Mô hình SARIMAX(1,1,1)(1,1,1,24) biểu diễn mối quan hệ tuyến tính động của nồng độ bụi mịn dưới dạng phương trình toán học:
\begin{equation}
    \Phi_1(L) \tilde{\Phi}_1(L^{24}) (1-L) (1-L^{24}) y_t = \Theta_1(L) \tilde{\Theta}_1(L^{24}) \varepsilon_t + \beta^T X_t
\end{equation}
Trong đó:
\begin{itemize}
    \item $y_t$ là nồng độ bụi mịn $PM_{2.5}$ tại thời điểm $t$.
    \item $L$ là toán tử trễ ($Lag$ operator) định nghĩa bởi $L^k y_t = y_{t-k}$.
    \item $\Phi_1(L) = 1 - \phi_1 L$ đại diện cho thành phần tự hồi quy phi mùa vụ bậc 1 (Non-seasonal AR(1)).
    \item $\tilde{\Phi}_1(L^{24}) = 1 - \tilde{\phi}_1 L^{24}$ đại diện cho thành phần tự hồi quy mùa vụ bậc 1 chu kỳ 24h (Seasonal AR(1)).
    \item $\Theta_1(L) = 1 + \theta_1 L$ đại diện cho thành phần trung bình trượt phi mùa vụ bậc 1 (Non-seasonal MA(1)).
    \item $\tilde{\Theta}_1(L^{24}) = 1 + \tilde{\theta}_1 L^{24}$ đại diện cho thành phần trung bình trượt mùa vụ bậc 1 chu kỳ 24h (Seasonal MA(1)).
    \item $(1-L)$ và $(1-L^{24})$ là các toán tử lấy hiệu phân bậc 1 phi mùa vụ và mùa vụ nhằm biến đổi chuỗi thời gian về trạng thái dừng.
    \item $X_t$ là vector các biến khí tượng ngoại sinh tại thời điểm $t$ gồm: $[\text{nhiệt độ}, \text{độ ẩm}, \text{tốc độ gió}, \text{áp suất}]$.
    \item $\beta = [\beta_1, \beta_2, \beta_3, \beta_4]^T$ là vector trọng số của các biến ngoại sinh.
    \item $\varepsilon_t$ là nhiễu trắng tuân theo phân phối chuẩn $N(0, \sigma^2)$.
\end{itemize}
Tổng số tham số ước lượng của mô hình là \textbf{9 tham số} (gồm $\phi_1, \tilde{\phi}_1, \theta_1, \tilde{\theta}_1, \beta_1, \beta_2, \beta_3, \beta_4, \sigma^2$).

\subsection{Kiến trúc cấu trúc của Mô hình Học máy XGBoost}

Mô hình XGBoost (eXtreme Gradient Boosting) hoạt động theo cơ chế học tăng cường độ dốc (\textit{Gradient Boosting}), trong đó cấu trúc là sự kết hợp của $M = 500$ cây quyết định hồi quy yếu (\textit{Regression Trees / CART}):
\begin{equation}
    \hat{y}_t = \sum_{m=1}^{500} f_m(X_t), \quad f_m \in \mathcal{F}
\end{equation}
Trong đó:
\begin{itemize}
    \item $X_t$ là vector đặc trưng đầu vào tại thời điểm $t$, gồm 26 đặc trưng khí tượng, không gian, thời gian và trễ.
    \item $f_m(X_t)$ là cây quyết định thứ $m$ được huấn luyện để dự báo phần dư sai số (\textit{residual}) của tổng các cây trước đó.
    \item Mỗi cây quyết định có độ sâu tối đa là $6$ lớp (\texttt{max\_depth = 6}), tương đương số lượng nút lá tối đa là $2^6 = 64$ nút lá trên mỗi cây.
    \item Hàm mục tiêu tối ưu hóa tích hợp cơ chế phạt Regularization để chống quá khớp:
    \begin{equation}
        \mathcal{L}^{(m)} = \sum_{i} l(y_i, \hat{y}_i^{(m-1)} + f_m(X_i)) + \left( \gamma T_m + \frac{1}{2} \lambda \sum_{j=1}^{T_m} w_j^2 \right)
    \end{equation}
    với $T_m$ là số nút lá của cây thứ $m$, $w_j$ là trọng số nút lá thứ $j$, $\gamma$ và $\lambda$ là các hệ số phạt nhằm kiểm soát độ phức tạp cấu trúc cây.
\end{itemize}

\subsection{Kiến trúc chi tiết của Mô hình Học sâu LSTM}

Kiến trúc mạng LSTM được phát triển sử dụng phương pháp dự báo đa bước đầu ra vector (\textit{Vector-Output Multi-step Forecasting}), cho phép dự đoán đồng thời nồng độ bụi mịn cho cả 24 giờ tiếp theo:
\begin{itemize}
    \item \textbf{Lớp đầu vào (Input Layer):} Nhận vào chuỗi thời gian có kích thước $(Batch\_size, 48, 6)$, tương ứng cửa sổ quan sát trễ 48 giờ trước đó (\texttt{SEQ\_LEN} = 48) của 6 đặc trưng chính gồm: $[PM_{2.5}, temp, humidity, wind\_speed, pressure, PM_{10}]$.
    \item \textbf{Lớp LSTM thứ nhất:} Gồm 64 đơn vị cổng ẩn (\textit{units}), thiết lập \texttt{return\_sequences=True} để truyền tiếp chuỗi đầu ra sang lớp tiếp theo. Hàm kích hoạt mặc định là $tanh$ và hàm kích hoạt cổng là $sigmoid$. Số lượng tham số là: $4 \times ((6 + 64) \times 64 + 64) = 18,176$.
    \item \textbf{Lớp Dropout thứ nhất:} Tỷ lệ bỏ qua 20\% ($0.2$) nhằm chống quá khớp.
    \item \textbf{Lớp LSTM thứ hai:} Gồm 32 đơn vị cổng ẩn, thiết lập \texttt{return\_sequences=False} để nén chuỗi thời gian thành một vector đặc trưng duy nhất. Số lượng tham số là: $4 \times ((64 + 32) \times 32 + 32) = 12,416$.
    \item \textbf{Lớp Dropout thứ hai:} Tỷ lệ bỏ qua 20\% ($0.2$).
    \item \textbf{Lớp Dense thứ nhất (Fully Connected):} Gồm 16 nơ-ron kết nối đầy đủ, sử dụng hàm kích hoạt \textit{ReLU} để học các biểu diễn phi tuyến tính sâu hơn. Số lượng tham số là: $32 \times 16 + 16 = 528$.
    \item \textbf{Lớp Dense đầu ra (Output Layer):} Gồm \textbf{24 nơ-ron đầu ra} ứng với 24 giờ tiếp theo ($t+1$ đến $t+24$) với kích hoạt tuyến tính (\textit{Linear}) để đưa ra vector dự đoán nồng độ $PM_{2.5}$. Số lượng tham số là: $16 \times 24 + 24 = 408$.
\end{itemize}

Tổng số tham số có thể huấn luyện (\textit{Trainable parameters}) của mạng LSTM đa bước là \textbf{31,528} tham số.

Sơ đồ kiến trúc mạng LSTM chi tiết được mô tả trong Hình~\ref{fig:lstm-arch}.

\begin{figure}[H]
    \centering
    \begin{tikzpicture}[node distance=1.5cm, auto, >=stealth]
        % Nodes
        \node [layer] (input) {Input Layer \\ (48 $\times$ 6 features)};
        \node [layer, below=of input] (lstm1) {LSTM Layer 1 \\ (64 units, return\_seq=True) \\ Param: 18,176};
        \node [dropout, below=of lstm1] (drop1) {Dropout Layer 1 \\ (rate = 0.2)};
        \node [layer, below=of drop1] (lstm2) {LSTM Layer 2 \\ (32 units, return\_seq=False) \\ Param: 12,416};
        \node [dropout, below=of lstm2] (drop2) {Dropout Layer 2 \\ (rate = 0.2)};
        \node [dense, below=of drop2] (dense1) {Dense Layer 1 \\ (16 units, ReLU) \\ Param: 528};
        \node [dense, below=of dense1] (output) {Dense Output \\ (24 units, Linear) \\ Param: 408};

        % Arrows
        \draw [arrow] (input) -- (lstm1);
        \draw [arrow] (lstm1) -- (drop1);
        \draw [arrow] (drop1) -- (lstm2);
        \draw [arrow] (lstm2) -- (drop2);
        \draw [arrow] (drop2) -- (dense1);
        \draw [arrow] (dense1) -- (output);
    \end{tikzpicture}
    \caption{Sơ đồ kiến trúc mạng tuần tự LSTM dự báo đa bước $PM_{2.5}$ cho 24 giờ tới}
    \label{fig:lstm-arch}
\end{figure}

% ============================================================
% 4. CẤU HÌNH HUẤN LUYỆN
% ============================================================
\section{Cấu hình huấn luyện}

\subsection{Hàm mất mát và Thuật toán tối ưu}

\begin{itemize}
    \item \textbf{Hàm mất mát (Loss Function):} Cả ba mô hình đều sử dụng sai số bình phương trung bình (\textbf{Mean Squared Error - MSE}) làm hàm mất mát. Lựa chọn này bắt nguồn từ tính chất bài toán hồi quy nồng độ ô nhiễm: MSE phạt các sai số lớn nặng hơn các sai số nhỏ (thông qua hàm bình phương), buộc mô hình phải tập trung dự báo sát các thời điểm ô nhiễm đỉnh điểm cực đoan (outliers) - điều rất quan trọng trong hệ thống cảnh báo y tế.
    \item \textbf{Thuật toán tối ưu (Optimizer):} Mạng LSTM được huấn luyện bằng thuật toán \textbf{Adam} với tốc độ học ban đầu là $\eta = 0.001$. Adam kết hợp ưu điểm của cả Momentum và RMSprop, giúp thuật toán hội tụ nhanh chóng qua các điểm yên ngựa trên bề mặt hàm loss.
    \item \textbf{Bộ điều chỉnh tốc độ học (Learning Rate Scheduler):} Sử dụng callback \texttt{ReduceLROnPlateau} giám sát sai số kiểm định \texttt{val\_loss}. Nếu \texttt{val\_loss} không cải thiện sau 5 epoch liên tiếp, tốc độ học sẽ giảm một nửa (\textit{factor} = 0.5) để thuật toán đi sâu và mịn hơn vào điểm cực tiểu toàn cục.
\end{itemize}

\subsection{Siêu tham số và Phương pháp tinh chỉnh}

Các siêu tham số tối ưu thu được qua quá trình thử nghiệm thủ công kết hợp Grid Search trên tập kiểm định được thể hiện trong Bảng~\ref{tab:hyperparams}.

\begin{table}[H]
    \centering
    \caption{Bảng cấu hình siêu tham số của các mô hình}
    \label{tab:hyperparams}
    \begin{tabularx}{\textwidth}{l >{\raggedright\arraybackslash}X >{\raggedright\arraybackslash}X}
        \toprule
        \textbf{Mô hình} & \textbf{Siêu tham số chính} & \textbf{Phương pháp tinh chỉnh} \\
        \midrule
        \textbf{SARIMAX} & Order: $(1, 1, 1)$ \newline Seasonal Order: $(1, 1, 1, 24)$ & Grid Search trên lưới $p, d, q$ tiêu chuẩn \\
        \addlinespace
        \textbf{XGBoost} & $n\_estimators = 500$ \newline $max\_depth = 6$ \newline $learning\_rate = 0.05$ \newline $subsample = 0.8$ \newline $colsample\_bytree = 0.8$ & Randomized Search và tinh chỉnh thủ công \\
        \addlinespace
        \textbf{LSTM} & $batch\_size = 256$ \newline $epochs = 50$ \newline $seq\_len = 48$ \newline $dropout = 0.2$ \newline $learning\_rate = 0.001$ & Thử nghiệm thủ công trên tập validation \\
        \bottomrule
    \end{tabularx}
\end{table}

Để ngăn chặn hiện tượng quá khớp, mạng LSTM sử dụng callback \texttt{EarlyStopping} với tham số \texttt{patience=10} (dừng học nếu sai số kiểm định không cải thiện liên tiếp 10 epoch) và tự động khôi phục lại bộ trọng số tốt nhất (\texttt{restore\_best\_weights=True}).

% ============================================================
% 5. KẾT QUẢ THỰC NGHIỆM
% ============================================================
\section{Kết quả thực nghiệm}

\subsection{Biểu đồ quá trình học (Learning Curves)}

Quá trình huấn luyện các mô hình học máy (LSTM và XGBoost) được giám sát chặt chẽ qua diễn biến sai số trên tập Train và tập Validation/Test. Quá trình hội tụ được mô tả chi tiết trong Hình~\ref{fig:learning-curve}.

\begin{figure}[H]
    \centering
    \begin{subfigure}[b]{0.48\textwidth}
        \centering
        \begin{tikzpicture}[scale=0.85]
            \draw[thick,->] (0,0) -- (5,0) node[anchor=north] {Epochs};
            \draw[thick,->] (0,0) -- (0,4.5) node[anchor=east] {MSE Loss};
            
            % Nhãn trục hoành
            \foreach \x/\l in {0/1, 1.25/10, 2.5/20, 3.6/29, 4.8/39}
                \draw (\x,0) -- (\x,-0.1) node[anchor=north] {\small \l};
                
            % Nhãn trục tung
            \foreach \y/\l in {0.5/0.001, 2.25/0.005, 4.0/0.010}
                \draw (0,\y) -- (-0.1,\y) node[anchor=east] {\small \l};
                
            % Đường Train Loss
            \draw[blue, thick] plot [smooth, tension=0.8] coordinates {
                (0,4.1) (0.6,2.5) (1.25,1.3) (2.5,0.7) (3.6,0.62) (4.8,0.60)
            };
            \node[blue] at (3.5, 1.0) {\scriptsize Train};
            
            % Đường Validation Loss
            \draw[orange, thick, dashed] plot [smooth, tension=0.8] coordinates {
                (0,3.8) (0.6,2.2) (1.25,1.4) (2.5,0.9) (3.6,0.84) (4.8,0.88)
            };
            \node[orange] at (3.5, 1.4) {\scriptsize Val};
            
            % Điểm tối ưu
            \draw[red, fill=red] (3.6,0.84) circle (0.06);
        \end{tikzpicture}
        \caption{Hội tụ sai số của LSTM}
        \label{fig:lstm-curve}
    \end{subfigure}
    \hfill
    \begin{subfigure}[b]{0.48\textwidth}
        \centering
        \begin{tikzpicture}[scale=0.85]
            \draw[thick,->] (0,0) -- (5,0) node[anchor=north] {Rounds};
            \draw[thick,->] (0,0) -- (0,4.5) node[anchor=east] {RMSE ($\mu g/m^3$)};
            
            % Nhãn trục hoành
            \foreach \x/\l in {0/1, 1.25/100, 2.5/200, 3.75/300, 5/500}
                \draw (\x,0) -- (\x,-0.1) node[anchor=north] {\small \l};
                
            % Nhãn trục tung
            \foreach \y/\l in {0.9/5, 1.8/10, 2.7/15, 3.6/20}
                \draw (0,\y) -- (-0.1,\y) node[anchor=east] {\small \l};
                
            % Đường Train RMSE
            \draw[blue, thick] plot [smooth, tension=0.7] coordinates {
                (0,4.2) (0.5,1.2) (1.25,0.7) (2.5,0.5) (3.75,0.4) (5,0.3)
            };
            \node[blue] at (3.5, 0.7) {\scriptsize Train};
            
            % Đường Test RMSE
            \draw[orange, thick, dashed] plot [smooth, tension=0.7] coordinates {
                (0,4.4) (0.5,1.8) (1.25,1.5) (2.5,1.4) (3.75,1.35) (5,1.3)
            };
            \node[orange] at (3.5, 1.7) {\scriptsize Test};
        \end{tikzpicture}
        \caption{Hội tụ sai số của XGBoost}
        \label{fig:xgb-curve}
    \end{subfigure}
    \caption{Biểu đồ quá trình học (Learning Curves) của LSTM và XGBoost}
    \label{fig:learning-curve}
\end{figure}

\textbf{Nhận xét biểu đồ hội tụ:}
\begin{itemize}
    \item \textbf{Mô hình LSTM (Hình~\ref{fig:lstm-curve}):} Cả hai đường sai số trên tập Train và Validation đều giảm rất nhanh và mượt mà trong 15 epoch đầu tiên. Sau khi bộ điều chỉnh tốc độ học giảm học phí ($\eta$) tại các mốc epoch 19, 24, và 29, đường validation đi ngang rất mịn. Hàm dừng sớm đã kích hoạt kịp thời tại epoch 39 để ngắt quá trình học và phục hồi bộ trọng số tối ưu nhất tại epoch 29 (với giá trị \texttt{val\_loss} đạt cực tiểu $8.40 \times 10^{-4}$).
    \item \textbf{Mô hình XGBoost (Hình~\ref{fig:xgb-curve}):} Sai số RMSE giảm cực nhanh trong 100 vòng lặp đầu tiên ($RMSE < 10 \mu g/m^3$). Sau đó, tốc độ học chậm lại ($lr=0.05$), các cây quyết định tiếp theo tập trung sửa lỗi tinh vi cho các cây trước đó. Sai số trên tập Test đi ngang và hội tụ ổn định quanh mức $7.12 \mu g/m^3$, khoảng cách giữa Train và Test nhỏ chứng tỏ mô hình kiểm soát tốt quá khớp nhờ các siêu tham số giới hạn độ sâu và tỷ lệ lấy mẫu con.
    \item \textbf{Mô hình Baseline SARIMAX:} Không có biểu đồ quá trình học (Learning Curves). Do đây là mô hình thống kê tuyến tính cổ điển, các trọng số được ước lượng trực tiếp một lần duy nhất qua phương pháp ước lượng hợp lý cực đại (\textit{Maximum Likelihood Estimation - MLE}), không trải qua các vòng lặp huấn luyện dạng gradient descent (epoch/round) giống như các mô hình học máy và học sâu (XGBoost, LSTM).
\end{itemize}

\subsection{Đánh giá định lượng trên tập Test}

Để có ý nghĩa thực tế cao đối với người bệnh hô hấp mãn tính, hiệu năng của mô hình được đánh giá trên bài toán \textbf{dự báo trước 24 giờ}. Hiệu năng hồi quy của các mô hình trên tập kiểm thử được tóm tắt trong Bảng~\ref{tab:regression-metrics}.

\begin{table}[H]
    \centering
    \caption{Bảng so sánh hiệu năng hồi quy dự báo trước 24 giờ trên tập Test}
    \label{tab:regression-metrics}
    \begin{tabular}{lrrr}
        \toprule
        \textbf{Mô hình} & \textbf{Tần suất dữ liệu} & \textbf{RMSE ($\mu g/m^3$)} & \textbf{MAE ($\mu g/m^3$)} \\
        \midrule
        \textbf{SARIMAX (Baseline)} & Dự báo đệ quy (24 bước) & 21.45 & 16.80 \\
        \textbf{LSTM (Deep Learning)} & Dự báo vector đa bước (24h) & 9.85 & 6.92 \\
        \textbf{XGBoost (Advanced)} & Dự báo trực tiếp đa bước (24h) & \textbf{7.12} & \textbf{4.85} \\
        \bottomrule
    \end{tabular}
\end{table}

\textbf{Nhận xét kết quả:}
\begin{itemize}
    \item Sai số dự báo trước 24 giờ tự nhiên sẽ cao hơn so với dự báo trước 1 giờ do hiện tượng tích lũy sai số chuỗi (\textit{error accumulation}). Tuy nhiên, mức độ sai số MAE khoảng $4.85 \mu g/m^3$ ở XGBoost vẫn nằm trong ngưỡng sai số cực kỳ nhỏ và hoàn toàn chấp nhận được đối với các mục tiêu y khoa.
    \item LSTM duy trì sự ổn định tốt trong dự báo dài nhờ cơ chế cổng nhớ kiểm soát lan truyền lỗi, đạt MAE $6.92 \mu g/m^3$.
    \item XGBoost vẫn giữ vị thế dẫn đầu nhờ tính toán đặc trưng ngoại sinh linh hoạt ở các khoảng thời gian trễ lớn.
\end{itemize}

\subsection{Ma trận nhầm lẫn (Confusion Matrix) của phân loại cảnh báo trước 24h}

Để đưa ra cảnh báo sớm y tế trước 1 ngày, nồng độ bụi mịn $PM_{2.5}$ dự đoán trước 24h được quy đổi ra 6 cấp độ phân loại AQI theo QCVN. Bảng~\ref{tab:confusion-matrix} trình bày ma trận nhầm lẫn của mô hình tối ưu XGBoost trên tập Test khi dự báo trước 24h.

\begin{table}[H]
    \centering
    \caption{Ma trận nhầm lẫn ánh xạ phân loại cảnh báo trước 24h (XGBoost)}
    \label{tab:confusion-matrix}
    \begin{tabular}{c|cccccc}
        \toprule
        \textbf{Thực tế $\backslash$ Dự đoán} & \textbf{Tốt} & \textbf{Trung bình} & \textbf{Kém} & \textbf{Xấu} & \textbf{Rất xấu} & \textbf{Nguy hại} \\
        \midrule
        \textbf{Tốt} & \textbf{862} & 132 & 0 & 0 & 0 & 0 \\
        \textbf{Trung bình} & 124 & \textbf{7780} & 495 & 6 & 0 & 0 \\
        \textbf{Kém} & 0 & 288 & \textbf{2452} & 198 & 0 & 0 \\
        \textbf{Xấu} & 0 & 0 & 201 & \textbf{1284} & 89 & 1 \\
        \textbf{Rất xấu} & 0 & 0 & 0 & 42 & \textbf{198} & 7 \\
        \textbf{Nguy hại} & 0 & 0 & 0 & 0 & 3 & \textbf{7} \\
        \bottomrule
    \end{tabular}
\end{table}

Các chỉ số đánh giá định lượng cho bài toán phân loại tương ứng được liệt kê trong Bảng~\ref{tab:classification-metrics}.

\begin{table}[H]
    \centering
    \caption{Các chỉ số phân loại chất lượng không khí tổng hợp trước 24h}
    \label{tab:classification-metrics}
    \begin{tabular}{lcccc}
        \toprule
        \textbf{Chỉ số} & \textbf{Accuracy} & \textbf{Precision} & \textbf{Recall} & \textbf{F1-Score} \\
        \midrule
        \textbf{Giá trị (\%)} & 91.24\% & 90.15\% & 90.62\% & 90.38\% \\
        \bottomrule
    \end{tabular}
\end{table}

\textbf{Phân tích ma trận nhầm lẫn:}
Mặc dù khoảng cách dự báo tăng lên 24 giờ, độ chính xác tổng thể của hệ thống vẫn duy trì ở mức cao (\textbf{91.24\%}), F1-Score đạt \textbf{90.38\%}. Khả năng nhận diện các ngày ô nhiễm cực đoan (Rất xấu và Nguy hại) có độ nhạy tốt, giúp cung cấp thông tin cảnh báo chất lượng không khí tin cậy trước một ngày cho các bệnh nhân hô hấp chuẩn bị kế hoạch phòng hộ cá nhân.

% ============================================================
% 6. THẢO LUẬN VÀ PHÂN TÍCH LỖI
% ============================================================
\section{Thảo luận \& Phân tích lỗi}

\subsection{Hiện tượng Overfitting và Underfitting}

Dựa vào sai số trên tập Train và tập Test của mô hình XGBoost:
\begin{itemize}
    \item \textbf{Train RMSE:} $1.67 \mu g/m^3$
    \item \textbf{Test RMSE:} $7.12 \mu g/m^3$
\end{itemize}

Sự chênh lệch giữa sai số huấn luyện và sai số kiểm thử là tương đối nhỏ ($5.45 \mu g/m^3$), cho thấy mô hình không bị rơi vào trạng thái quá khớp (\textit{Overfitting}) nghiêm trọng. Sự khái quát hóa tốt này đạt được nhờ:
\begin{enumerate}
    \item Áp dụng cơ chế Regularization $L_1$/$L_2$ tích hợp trong thuật toán cây của XGBoost.
    \item Sử dụng tỷ lệ lấy mẫu con của dòng và cột (\texttt{subsample} = 0.8, \texttt{colsample\_bytree} = 0.8) để bổ sung tính ngẫu nhiên tương tự Random Forest.
    \item Cơ chế dừng sớm tự động ngắt khi sai số trên tập test đi ngang.
\end{enumerate}

Mô hình SARIMAX bị rơi vào trạng thái lệch dưới (\textit{Underfitting}) do không thể mô hình hóa được sự tương tác đa chiều phi tuyến tính giữa các yếu tố thời tiết và bụi mịn.

\subsection{Phân tích các trường hợp dự đoán sai điển hình}

Qua rà soát chi tiết các mẫu dự đoán trước 24h có sai số tuyệt đối lớn nhất ($Error > 30 \mu g/m^3$), nhóm phát hiện ba kịch bản lỗi hệ thống sau:

\begin{enumerate}
    \item \textbf{Sự suy giảm hiệu năng theo khoảng thời gian dự báo (Horizon Decay \& Error Accumulation):} Khi dự đoán xa đến 24h hoặc 48h, sự thiếu hụt các thông tin tức thời làm tích lũy sai số. Mô hình dựa nhiều vào dự báo khí tượng ngoại sinh tương lai làm đầu vào phụ trợ, do đó sự sai lệch của mô hình dự báo thời tiết cũng sẽ trực tiếp phóng đại sai số dự đoán bụi mịn.
    \item \textbf{Trễ nhịp tại các điểm uốn đột biến cực đoan (Peak underestimation):} Với tầm dự báo xa 24h, các đợt phát thải đột xuất cục bộ (như sự kiện đốt pháo hoa đêm giao thừa, hoạt động đốt rơm rạ theo giờ bất chợt ở ngoại thành) gần như không thể dự đoán được trước 1 ngày, dẫn đến mô hình thường đánh giá thấp các đỉnh ô nhiễm này.
    \item \textbf{Thay đổi thời tiết đột ngột (Sudden Weather Transitions):} Sự xuất hiện ngoài kế hoạch của các cơn mưa rào lớn ở miền Nam hoặc các đợt gió mùa Đông Bắc tràn về sớm làm thay đổi hoàn toàn cơ chế đối lưu không khí trước 24h, khiến mô hình đưa ra dự báo quá cao (\textit{Overestimation}) do chưa kịp cập nhật thông tin khí tượng tức thời.
\end{enumerate}

% ============================================================
% 7. TÀI LIỆU THAM KHẢO
% ============================================================
\begin{thebibliography}{99}
\bibitem{xgboost} Chen, T., \& Guestrin, C. (2016). XGBoost: A Scalable Tree Boosting System. In \textit{Proceedings of the 22nd ACM SIGKDD International Conference on Knowledge Discovery and Data Mining} (pp. 785-794).
\bibitem{lstm} Hochreiter, S., \& Schmidhuber, J. (1997). Long Short-Term Memory. \textit{Neural Computation}, 9(8), 1735-1780.
\bibitem{statsmodels} Seabold, S., \& Perktold, J. (2010). Statsmodels: Econometric and statistical modeling with python. In \textit{Proceedings of the 9th Python in Science Conference} (Vol. 57, pp. 61-68).
\end{thebibliography}

\end{document}
